% Source-derived chapter 04: Curves of canonical  type
% Original source: no-enriques-surface-over-the-integers
\chapter{Curves of canonical  type}
\label{no-enriques-surface-over-the-integers-chapter-04}
\source{no-enriques-surface-over-the-integers}

\begin{remark}[Source section]
\label{no-enriques-surface-over-the-integers-chapter-04-source}
\source{no-enriques-surface-over-the-integers}
\source{no-enriques-surface-over-the-integers}
This chapter reproduces the corresponding section of the retrieved
LaTeX source, including its definitions, statements, examples, and
proofs. It has not yet been translated into Lean.
\notready
\end{remark}

% The source section title was: Curves of canonical  type
\mylabel{Curves}

Let $k$ be a ground field of characteristic $p\geq 0$, and $S$ be a smooth proper surface with $h^0(\O_S)=1$.
In this section we apply Theorem \ref{geometric and schematic fiber} to fibers in genus-one fibrations,
and collect additional facts for the situation that $k=\FF_q$ is finite.

We start with  some useful general terminology:
Let $C\subset S$ be a curve, and  $C=m_0C_0+\ldots+m_rC_r$ be the decomposition into irreducible components, together with
their multiplicities.
We say that $C$ is of \emph{fiber type} if  
$(C\cdot C_i)= 0$ for all indices  $0\leq i\leq r$.
Obviously, this condition transfers to multiples and connected components of the curve.
If $C$ is connected and $\gcd(m_0,\ldots,m_r)=1$, we say that
$C$ is an \emph{indecomposable curve of fiber type}. 
Then the intersection matrix $N=(C_i\cdot C_j)_{1\leq i,j\leq r}$ is negative-semidefinite,
and   the radical of
the intersection form on the lattice $L=\bigoplus_{i=0}^r \ZZ C_i$ is generated by $m_0C_0+\ldots+m_rC_r$.

Given a connected curve  of fiber type $C=\sum m_iC_i$, we call $m=\gcd(m_0,\ldots,m_r)$ its \emph{multiplicity}.
We say that $C$ is \emph{simple} if $m=1$, and \emph{multiple} otherwise.
Setting $n_i=m_i/m$, we call  $C_\ind=\sum n_iC_i$ the \emph{underlying indecomposable curve of fiber type}.

Suppose now that we have a contraction $f:S\ra B$ of fiber type onto some curve $B$.
For each geometric point $\ba:\Spec(\Omega)\ra B$, the geometric fiber $S_\ba=f^{-1}(\ba)$ is a connected 
curve of fiber type on the base-change $S\otimes_k\Omega$. It is useful to write $f^\inv_\ind(\ba)$ for the underlying
indecomposable curve of fiber type, and $f^\inv_\red(\ba)$ for its reduction.
An analogous notation will be used for  the schematic fiber $S_a=f^\inv(a)$
over the image point $a\in B$.
 
We are mainly interested in \emph{curves of canonical type}. By definition, this means 
$(C\cdot C_i)=(K_S\cdot C_i)=0$ for all indices $0\leq i\leq r$. The above locutions for curves of fiber type
are used in an analogous way for  curves of canonical types.

Suppose now that $f:S\ra B$ is a genus-one fibration that is relatively minimal.
Then every geometric fiber $ f^{-1}(\ba)$ is a connected curve of canonical type.
The underlying indecomposable curves of canonical type $f^\inv_\ind(\ba)$ were classified by N\'eron \cite{Neron 1964} and Kodaira \cite{Kodaira 1963}.
Throughout, we denote the \emph{Kodaira symbols} 
\begin{equation}
\label{kodaira symbols}
\source{no-enriques-surface-over-the-integers}
\I_m,\;\I_n^*,\;\II,\;\III,\;\IV,\;\IV^*,\;\III^*,\; \II^*
\end{equation}
to designate the geometric fibers.  The symbols $\I_m$, $m\geq 1$ are also called \emph{semistable} or \emph{multiplicative},
whereas $\I_n^*,\II, \ldots,\II^*$,  are referred to as \emph{unstable} or \emph{additive}. Here $n\geq 0$.
As customary, the symbol $\I_0$ indicates that $f^\inv_\ind(\ba)$ is an elliptic curve; we then distinguish the cases
that the  smooth curve is \emph{ordinary} or \emph{supersingular}.
By abuse of terminology 
we say that a geometric fiber is \emph{reducible or singular} if the scheme $f^\inv_\ind(\ba)$ has the respective property.
If this curve is  singular with  $f$  elliptic, or reducible with $f$ quasielliptic, or multiple, 
 we say that the geometric fiber is \emph{degenerate}.

The   classification of geometric fibers $f^\inv(\ba)$ was extended to schematic fibers $ f^{-1}(a)$ by Szydlo \cite{Szydlo 2004},
where the situation becomes considerably more involved.
If all irreducible components   of $S_a=f^\inv(a)$ 
are birational to $\PP^1$, the map $\Gamma(S_\ba)\ra\Gamma(S_a)$ is a graph isomorphism respecting edge labels,
and we also use the Kodaira symbols in \eqref{kodaira symbols} to designate   $f^{-1}_\ind(a)=m_1C_1+\ldots+m_rC_r$.
We need, however, the following modifications for $r\leq 2$, which play an important role in our later applications:

For $r=1$ the curve  $C=C_1$ has $h^0(\O_C)=h^1(\O_C)=1$, and the normalization $\nu:\PP^1\ra C$ is described by 
a  \emph{conductor square}
$$
\begin{CD}
\source{no-enriques-surface-over-the-integers}
\notready
\Spec(R)	@>>>	\PP^1\\
@VVV		@VV\nu V\\
\Spec(k)	@>>>	C,
\end{CD}
$$
which is both cartesian and cocartesian, see \cite{Fanelli; Schroeer 2020}, Appendix A for more details. 
Geometrically speaking, the closed subscheme $\Spec(R)\subset\PP^1$
is contracted to a rational point $\Spec(k)\subset C$.
Here $R$ is a finite $k$-algebra of length two. If $R=k\times k$ or $R=k[\epsilon]$ with $\epsilon^2=0$ 
then $C$ is the rational nodal curve or the rational cuspidal curve and we use the standard Kodaira symbol $\I_1$ and $\II$,
respectively. If  $R=E$ is a field, then $C$ is a  twisted form, 
and we use   \emph{twisted Kodaira symbols} $\tI_1$ and $\widetilde{\II}$ instead.
Note that in the twisted situation, a non-rational point on $\PP^1$ gets identified
with a rational point on $C$, with consequences for the size of $C(k)$.

A similar situation appears for $r=2$: Then   $C=C_1+C_2$ sits in the conductor square
$$
\begin{CD}
\source{no-enriques-surface-over-the-integers}
\notready
\Spec(R\times R)	@>>>	\PP^1\amalg\PP^1\\
@VVV		@VV\nu V\\
\Spec(R)	@>>>	C,
\end{CD}
$$
where $R$ is a $k$-algebra of length two. For $R=k\times k$ or $R=k[\epsilon]$, where $\epsilon^2=0$, we again use the standard Kodaira symbols
$\I_2$ and $\III$, respectively. If $R=E$ is a field extension, we have twisted forms and  use $\tI_2$ and $\widetilde{\III}$ instead.
Again, the twisted situation has   consequence for the size of $C(k)$.

Note that over perfect ground fields, and in particular over finite ground fields, for $r\leq 2$ the cases $\widetilde{\II}$ and $\widetilde{\III}$ do not occur,
and   we only have $\tI_1$ and $\tI_2$ as twisted Kodaira symbols. 
One then says that $C$ is a \emph{non-split semistable fiber}.  
For $r\geq 3$ there are further   possible twists, but then  $\Gamma(S_\ba)\ra\Gamma(S_a)$ is not a graph isomorphism.
We   record the following consequence of Theorem \ref{geometric and schematic fiber}:

\begin{proposition}
\source{no-enriques-surface-over-the-integers}
\notready
\mylabel{constant num over finite field}
Suppose the ground field $k=\FF_q$ is finite and
that  the  group scheme $\Num_{S/k}$ is constant.
Then $\Gamma(S_\ba)\ra\Gamma(S_a)$ is a graph isomorphism respecting edge labels, and 
the implications (i)$\Rightarrow$(ii)$\Leftrightarrow$(iii) hold  among the following three conditions:
\begin{enumerate}
\item The geometric fiber $S_\ba$ is reducible.
\item The closed point $a\in B$ is a rational point.
\item The irreducible components of the schematic fiber $S_a$ are birational to $\PP^1$.
\end{enumerate}
\end{proposition}

\proof
The   assertion on $\Gamma(S_\ba)\ra\Gamma(S_a)$ immediately follows from Theorem \ref{geometric and schematic fiber},
which also ensures that   $k=\kappa(a)$ provided that $S_a$ is reducible.
This already gives (i)$\Rightarrow$(ii). The implication (ii)$\Leftarrow$(iii) is trivial.
For the converse,  let $C\subset S_a$ be an irreducible component, and $C'\ra C$ be its normalization.
Then $C\otimes\Omega$ is reduced, with normalization $C'\otimes\Omega$, because the field extension $k\subset\Omega$
must be separable. The curve $C\otimes\Omega$ is  birational to $\PP^1_\Omega$, because 
$\Gamma(S_\ba)\ra\Gamma(S_a)$ is a graph isomorphism. Hence the same holds for $C'\otimes\Omega$.
In other words, $C'$ is a one-dimensional Brauer--Severi variety.  
By Wedderburn's Theorem, we must have $C'\simeq \PP^1$.
\qed

\begin{proposition}
\source{no-enriques-surface-over-the-integers}
\notready
\mylabel{canonical type rational points}
Assumptions as in Proposition \ref{constant num over finite field}.
Let $C=f^\inv(a)$ be a sche\-matic  fiber over a rational point $a\in\PP^1$, and $r\geq 1$ be the number
of irreducible components. Suppose the reduced scheme $C_\red$ is singular. Then the number  of rational points in the fiber 
is  given by the following table:
$$
\begin{array}[b]{@{}lllll}
\toprule
\text{\rm Kodaira symbol}	& \I_r	& \tI_1	&\tI_2	& \text{\rm unstable}\\
\midrule
\Card C(\FF_q)				& rq	& q+2 	& 2q+2	& rq+1 \\
\bottomrule
\end{array}
$$
For the field $k=\FF_2$ with two elements, the  Kodaira symbol $\I_0^*$ is impossible.
\end{proposition}

\proof
The computation of $n\geq 0$ follows directly by comparing $C=f^\inv_\red(a)$ with its normalization $C'=\PP^1\cup\ldots\cup\PP^1$, 
and is left to the  reader. If the Kodaira symbol is $\I_0^*$, then in $C=C_0+C_1+\ldots+C_4$
the terminal components $C_1,\ldots,C_4$ intersect the central component $C_0=\PP^1$ in four rational points.
This implies $4\leq \Card\PP^1(\FF_q)=q+1$, which means $q\neq 2$.
\qed

\medskip
We now   examine the situation that $k=\FF_2$  and  the scheme $ f^\inv_\ind(a)$ is smooth. 
This genus-one curve contains a rational point (\cite{Lang 1956}, Theorem 2), hence
can be regarded as an elliptic curve $E=f^\inv_\ind(a)$. 
Let us recall from \cite{Husemoller 1987}, Chapter 3, \S6 
that up to isomorphism, there are five elliptic curves $E_1,\ldots,E_5$ over the prime field $k=\FF_2$.
The groups $E_i(k)$ are cyclic, and all   $1\leq n\leq 5$ occur as orders. Deviating from loc.\ cit., we
choose indices in a canonical way according to the order of the group of rational points, and record:

\begin{proposition}
\source{no-enriques-surface-over-the-integers}
\notready
\mylabel{elliptic curves}
Up to isomorphisms, the elliptic curves $E_i$ over the field $k=\FF_2$ with $\Card E_i(\FF_2)=i$
are given by the following table:
$$
\begin{array}[b]{@{}lll}
\toprule
\text{elliptic curve}	& \text{Weierstra\ss{} equation}	& \text{invariant}\\
\toprule
E_1			& y^2+y=x^3+x^2+1 			& \text{supersingular with $j=0$}\\
E_3			& y^2+y=x^3\\
E_5			& y^2+y=x^3+x^2\\
\midrule
E_2			& y^2+xy=x^3+x^2+x  			& \text{ordinary with $j=1$}\\
E_4			& y^2+xy=x^3+x\\
\bottomrule
\end{array}
$$
\end{proposition}
 

%===========================================================
